% !PRINTED 123
\weberpage{123}
%<<CONT
is. Now if $c$ is an element of $P$ not belonging to $A$, then $x^{-1}cx$ will not be equal to $c$ for all elements $x$. But certain elements $b, b', \ldots$ will be found among the $x$, among them certainly the unit element, which satisfy the condition
\[ x^{-1}cx = c \]
and these elements form a group, because from
\weq{3}{b^{-1}cb = c,\quad b'^{-1}cb' = c}
it follows that also
\[ (b'b)^{-1}c(b'b) = b^{-1}b'^{-1}cb'b = c \]
must hold. We denote this group by $B$ and its degree and index by $\mu, \varepsilon$, so that
\weq{4}{n = \mu\varepsilon}
holds; $B$ is a proper divisor of $P$, and hence $\varepsilon$ is in any case greater than 1, since otherwise $c$ would be contained in $A$, contrary to the assumption.

If $g$ is an element of $P$ that does not occur in $B$, and $b$ an arbitrary element of $B$, then by (3):
\weq{5}{g^{-1}b^{-1}cbg = g^{-1}cg,}
which is different from $c$, because $g$ does not occur in $B$. If therefore one decomposes $P$ into the subsidiary groups:
\weq{6}{P = B + Bg_1 + Bg_2 + \cdots + Bg_{\varepsilon-1},}
then, on putting for $x$ in $x^{-1}cx$ all the elements of one of these subsidiary groups, one always obtains the same element $g^{-1}cg$, and one thus obtains $\varepsilon$ and no more elements, and each equally often:
\weq{7}{c,\ g_1^{-1}cg_1,\ g_2^{-1}cg_2,\ \ldots,\ g_{\varepsilon-1}^{-1}cg_{\varepsilon-1};}
all these elements, whose totality we shall denote by $C$, are different from one another and from the elements $a$. But $C$ is not a group, since the unit element does not occur among them.

If now the totality of the elements of $P$ is not yet exhausted by $A$ and $C$, one takes an element $c'$ occurring neither in $A$ nor in $C$ and forms by the same rule the group $B'$ of degree $\mu'$ of all elements which satisfy the condition
\[ b'^{-1}c'b' = c' \]
and the system $C'$:
\weq{8}{c',\ g_1'^{-1}c'g_1',\ g_2'^{-1}c'g_2',\ \ldots,\ g_{\varepsilon'-1}'^{-1}c'g_{\varepsilon'-1}',}
%CONT>>
